\documentclass[10pt,twocolumn]{article}
\usepackage[T1]{fontenc}
\usepackage[utf8]{inputenc}
\usepackage{lmodern}
\usepackage[margin=1.8cm,top=2.4cm,bottom=2cm,columnsep=0.6cm]{geometry}
\usepackage{fancyhdr}
\usepackage{titlesec}
\usepackage{booktabs}
\usepackage{amsmath,amssymb,amsfonts}
\usepackage{microtype}
\usepackage{xcolor}
\usepackage{mdframed}
\usepackage{parskip}
\usepackage{enumitem}
\usepackage{hyperref}

\definecolor{rulered}{RGB}{180,30,30}
\definecolor{boxgray}{RGB}{245,245,245}
\definecolor{keywordgray}{RGB}{100,100,100}

\pagestyle{fancy}
\fancyhf{}
\fancyhead[L]{\textsc{\small Claude Mag}}
\fancyhead[C]{\small\textit{Machine Learning Research Digest}}
\fancyhead[R]{\small 2026-09-14}
\fancyfoot[C]{\small\thepage}
\renewcommand{\headrulewidth}{0.8pt}

\titleformat{\section}{\normalsize\bfseries\color{rulered}}{}{0pt}{}%
  [\vspace{-4pt}\color{rulered}\rule{\columnwidth}{0.4pt}]
\titlespacing{\section}{0pt}{8pt}{4pt}

\hypersetup{colorlinks=true,urlcolor=rulered,linkcolor=black}

\begin{document}
\setlength{\parindent}{0pt}
\setlength{\parskip}{4pt}

{\small\color{keywordgray}\textbf{Keywords:}
  long video understanding \textbullet{}
  agentic AI \textbullet{}
  edge-cloud \textbullet{}
  visual routing \textbullet{}
  multimodal LLM}\par\vspace{4pt}

{\LARGE\bfseries\raggedright
  Caption-once, Frames-on-Demand:
  Visual-Need Routing for Budget-Aware
  Agentic Long Video Understanding}\par\vspace{4pt}

{\large\itshape\raggedright
  One Offline Captioning Pass, a Dual-Track Narrative Index,
  and a Lightweight Router Cut Frame API Calls by 71\%
  While Improving Accuracy}\par\vspace{6pt}

{\small\textbf{By} Weitong Cai, Hang Zhang, Yukai Huang,
  Yiqiao Xie, Shan Gao, Jiankang Deng, Songcen Xu,
  Jifei Song, Zhensong Zhang
  (Imperial College London; Huawei Noah's Ark Lab; et al.)}\par\vspace{2pt}

{\small\color{keywordgray}arXiv:2609.11899 \textbullet{} EMNLP 2026}
\par\vspace{2pt}\noindent{\color{rulered}\rule{\columnwidth}{0.6pt}}\vspace{6pt}

Deploying long-video understanding on edge devices imposes two
simultaneous constraints: tight compute budgets (no dense frame encoding
over hours of content) and tight bandwidth budgets (no per-query frame
refetch from the cloud). Prior systems sacrifice accuracy for
efficiency (sparse frame methods) or efficiency for accuracy (dense
MLLM methods). Caption-once, Frames-on-Demand (CFD) resolves the
tradeoff: a single offline captioning pass builds a reusable dual-track
narrative index, and a lightweight Visual-Need Router at query time
triggers pixel retrieval only when the question requires it --- 29\%
of the time.

\begin{mdframed}[backgroundcolor=boxgray,linecolor=rulered,linewidth=1pt,
    innerleftmargin=6pt,innerrightmargin=6pt,
    innertopmargin=6pt,innerbottommargin=6pt]
\textbf{\small AT A GLANCE}\par\vspace{4pt}
\begin{description}[leftmargin=0pt,labelsep=4pt,itemsep=2pt,parsep=0pt,
    font=\small\bfseries]
  \item[Problem:] Long-video understanding on edge devices is blocked
    by the cost of dense frame processing and per-query frame retrieval.
  \item[Why it matters:] Hours-long video QA is needed for wearables,
    in-car systems, and surveillance; bandwidth and compute are limited.
  \item[Method:] Caption video once; build event-level story skeleton
    + clip-level micro-log; route each query through Visual-Need Router;
    retrieve frames only for perceptual sub-questions.
  \item[Result:] 74.8\%/60.9\%/66.7\% on EgoSchema/LVBench/Video-MME-Long;
    average 1.16 frames per query vs.\ 18--64 for baselines.
  \item[Evidence:] Ablations on dual-track index, router design,
    and frame budget $R_{\max}$.
\end{description}
\end{mdframed}

\section{The Visual-Textual Duality}

The core insight is that long-video questions separate cleanly into
two types:

\textbf{Structural questions} (temporal order, causality, scene
transitions, narrative arc) are answered better from language memories
than from dense frames. A prose description of a sequence of events
conveys order and causality more compactly and reliably than 3,600
frames at 1 fps.

\textbf{Perceptual questions} (clothing color, on-screen text, object
state, character identity) require pixels. Language captions inevitably
abstract away discriminating visual details; the answer ``the person
is wearing a blue jacket'' is useless if the question distinguishes
between navy, teal, and cobalt.

Prior systems ignore this duality. Dense-frame MLLMs pay the cost of
all frames for all questions. Text-only memory agents fail systematically
on perceptual questions. CFD exploits the duality: build a reusable
language index and retrieve frames only when pixels are genuinely needed.

\section{The CFD Framework}

\textbf{Offline phase (run once at ingestion):}

\textbf{1. Clip segmentation.} Video is divided into clips
$\{c_1, \ldots, c_K\}$ at scene boundaries by a local shot-detection
model (no API calls).

\textbf{2. Clip captioning.} Each clip $c_k$ is captioned by a
compact 3B-parameter local MLLM, producing micro-log entry $m_k$:
timestamp range, key actors/objects, described actions, verbatim
on-screen text.

\textbf{3. Story synthesis.} A cloud MLLM writes an event-level story
skeleton $S$ (200--500 words) covering the main narrative arc, key
events, and causal relationships from the full micro-log sequence.

\textbf{4. Index construction.} The dual-track index
$\mathcal{I} = (S,\, \{m_k\},\, \{t_k\})$ is stored locally.

\textbf{Online phase (per query):}

\textbf{1. Story pass.} Read $S$ and identify candidate time windows
$\mathcal{W}^*$ relevant to the query.

\textbf{2. Micro-log pass.} Read micro-logs for $\mathcal{W}^*$ and
attempt a text-only answer.

\textbf{3. Visual-Need Router.} Classify remaining uncertainty as
``perceptual'' or ``structural''. If structural: return text answer.
If perceptual: request bounded keyframe retrieval.

\textbf{4. Bounded retrieval.} Request at most $R_{\max} = 4$
frames from the relevant window; integrate into the final response.

\section{Technical Formulation}

The full inference pipeline computes:
\[
  \hat{\mathcal{A}} = \mathrm{MLLM}\!\left(
    q,\; S,\; \{m_k : t_k \in \mathcal{W}^*\},\;
    \{f_j : j \in \mathcal{R}(q)\}
  \right)
\]
where the frame retrieval set is:
\[
  \mathcal{R}(q) = \begin{cases}
    \arg\max_{|J| \leq R_{\max}} \mathrm{Rel}(q, \{f_j\}_{j \in J})
      & \phi(q, \mathcal{I}) = \mathrm{perceptual} \\
    \emptyset & \text{otherwise}
  \end{cases}
\]
and $\phi: (q, \mathcal{I}) \to \{\mathrm{perceptual}, \mathrm{structural}\}$
is the Visual-Need Router.

\textbf{Expected cost per query:}
\[
  \mathbb{E}[|\mathcal{R}(q)|] = R_{\max} \cdot p_{\mathrm{perc}}
  = 4 \times 0.29 = 1.16 \text{ frames}
\]
vs.\ $\sim$64 for dense-frame baselines and $\sim$18--22 for prior
agentic methods.

\section{Visual-Need Router}

The router $\phi$ is a 350M-parameter BERT-style encoder trained on
2,400 (question, micro-log) pairs annotated for visual necessity. It
classifies a query as ``perceptual'' if the discriminating visual
attribute cannot be recovered from the narrative text alone.

Router accuracy on 600 held-out pairs: 91.4\%. Inference latency:
${<}1$\,ms on CPU --- negligible vs.\ frame retrieval latency
(200--800\,ms round-trip).

\section{Experimental Results}

\begin{center}
\begin{tabular}{lccc}
\toprule
Method & EgoS. & LVB. & VME-L \\
\midrule
GPT-4V-Long (dense)  & 72.4\% & 58.3\% & 64.1\% \\
VideoAgent           & 68.9\% & 52.7\% & 59.8\% \\
EgoVideoQA           & 70.3\% & 54.1\% & 61.2\% \\
\textbf{CFD (ours)}  & \textbf{74.8\%} & \textbf{60.9\%} & \textbf{66.7\%} \\
\bottomrule
\end{tabular}
\end{center}

\vspace{2pt}
{\small EgoS. = EgoSchema; LVB. = LVBench; VME-L = Video-MME Long.}
\vspace{4pt}

CFD improves over the best prior method by 2.4--2.8 accuracy points
while using 15--55$\times$ fewer frames per query.

\section{Ablations and Interpretation}

\textbf{Without story skeleton:} Accuracy drops 5.3 points on LVBench.
Multi-scene temporal reasoning requires the sketch-level narrative.

\textbf{Without micro-logs:} Accuracy drops 3.1 points on EgoSchema.
Clip-level detail is needed for fine-grained activity understanding.

\textbf{Without Visual-Need Router (always retrieve):} Accuracy
improves by only 0.9 points but frame cost increases 3.4$\times$.
The router's accuracy--efficiency tradeoff is strongly favorable.

\textbf{Single-track index:} Either track alone scores 4--8 points
below the dual-track system, confirming the two tracks are complementary.

\section*{Reference}

Weitong Cai et al.\ \textbf{Caption-once, Frames-on-Demand: Visual-Need
Routing for Budget-Aware Agentic Long Video Understanding.}
arXiv:2609.11899, EMNLP 2026.
\url{https://arxiv.org/abs/2609.11899}

\end{document}
